\section{Conclusions}
This study presented DNA+, a bio-inspired requirements engineering framework that addresses the persistent representational and operational limitations of automated software requirements classification. By translating natural language specifications into synthetic biological genomic sequences through a dynamic syntactic-semantic dinucleotide codon grammar, DNA+ resolves the contextual decoupling inherent to prior static single-nucleotide mapping techniques. The formulation constructs a compact 23-dimensional alignment-free physicochemical profile—incorporating sixteen dinucleotide relative abundance ratios, GC-content, AT/GC skew, and Shannon sequence entropy—that provides an invariant structural signature for each requirement statement. Fusing these genomic properties with calibrated lexical frequencies and dense Sentence-BERT contextual embeddings establishes a 457-dimensional multimodal manifold that captures high-order semantic abstractions alongside low-level grammatical dependencies.

Comprehensive empirical evaluation across three diverse benchmark corpora comprising 8,654 authentic requirements (canonical PROMISE, extended PROMISE-EXP, and industrial FNFC) demonstrated the decisive superiority of the proposed framework. The champion Gradient Boosting and Decision Tree architectures achieved peak classification accuracies of 94.77\% and 94.86\% with macro-F1 scores of 85.61\% and 88.54\% on the industrial benchmark, alongside 90.40\% accuracy on the extended multi-class corpus under severe class imbalance (37.0:1). Rigorous parametric and non-parametric statistical hypothesis testing confirmed that observed performance improvements over baseline representations are statistically significant ($t = 11.6865, p < 0.001$; Wilcoxon $W = 28.0, p < 0.01$) with massive standardized effect sizes (Cohen's $d = 4.4171$, Cliff's $\delta = 1.0000$). Furthermore, direct comparative benchmarking against five premier 2025--2026 state-of-the-art architectures demonstrated a sustained accuracy lead ranging from +5.17\% to +9.67\% without requiring synthetic data oversampling or costly GPU fine-tuning.

Beyond raw predictive accuracy, the integrated dual-tier explainable artificial intelligence framework combines global TreeSHAP game-theoretic attribution with local LIME token-level justifications, providing software engineering practitioners with verified interpretability and auditability. By establishing that alignment-free biological sequence physics inherently stabilize natural language classification under extreme distribution skew, this investigation provides a scalable, computationally efficient paradigm for automated quality attribute triage in continuous software development lifecycles.

Future research will explore three primary directions. First, expanding the dynamic dinucleotide codon grammar to tri-nucleotide reading frames ($k=3$, 64-codon space) to capture complex multi-clause contractual constraints and nested architectural dependencies. Second, evaluating cross-lingual transferability across non-English and mixed-language industrial requirements repositories. Third, integrating the DNA+ lightweight feature extraction pipeline directly into commercial requirements management platforms and automated test-case generation environments for real-time specification quality assurance.
